We trained DGCNN models on the BNCI 2014-001 Motor Imagery dataset and analysed the resulting learned graph representations. The primary objective was to answer our research question of whether the learned graph representations exhibit any form of underlying stability across runs. For this purpose we employed both model similarity analysis using CKA and a range of graph centrality measures. 
\vspace{0.5em}\\
Our findings reveal a pattern: the electrodes C3 and CP3 consistently emerged as the central nodes across several centrality metrics, namely barycentres, betweenness centrality, Laplacian centrality, and current-flow betweenness centrality. Thus, for our particular experiment, we observe a certain degree of consistency in which electrodes have high importance scores in the learned graphs according to the aforementioned centrality measures.
\vspace{0.5em}\\
Given that these electrodes are located over the motor cortex, area involved in motor planning and execution, this finding suggests that the models are learning structured graph representations, and are also capturing meaningful and interpretable neural features specific to the motor imagery task.
\vspace{0.5em}\\
While these results are promising, the study was limited to a single dataset and a single GNN architecture. Broader validation across multiple EEG datasets and diverse neural network configurations is necessary to determine the generalizability of these findings. Nevertheless, the observed emergence of biologically related graph structures highlights the potential of GNNs for learning interpretable and stable representations in EEG analysis, providing motivation for future research in this area.